В настоящем документе были продемонстрированы основные возможности шаблона \texttt{faimatantemplate}: семантические блоки, общая нумерация, именованные ссылки, дополнительный материал, формулы, рисунки, таблицы, алгоритмы, программный код и задачи со скрываемыми решениями.
При создании нового проекта рекомендуется сохранить структуру стилевых файлов и сначала определить общие соглашения о математических обозначениях, метках и оформлении ссылок.
Новые команды следует добавлять в файл \verb|style/helpers.tex| и одновременно документировать в соответствующем разделе руководства по стилю.
Такой порядок позволяет поддерживать единообразие документа и упрощает совместную работу нескольких авторов.
Шаблон можно свободно изменять и использовать как основу для собственных проектов при условии указания автора дизайна и \LaTeX-оформления в выходных данных итогового документа.